\documentclass[10pt,a4paper]{amsart}
\usepackage[margin=19mm]{geometry}
\usepackage{amsmath,amssymb,mathtools}
\usepackage[scr=rsfs,cal=euler]{mathalfa}
\usepackage{xfrac}
\usepackage[unicode,hidelinks,bookmarks=false]{hyperref}
\newcommand{\ie}{i.e.\ }
\newcommand{\cf}{cf.\ }
\newcommand{\resp}{respectively\ }
\newcommand{\Subsection}{\subsection}
\providecommand{\nobreakdash}{\nobreak\hbox{-}}
\setlength{\emergencystretch}{2em}
\multlinegap0pt
\usepackage{tabularx,booktabs}
\usepackage{fullpage}
\usepackage{lipsum}
\usepackage{bm}
\usepackage{algorithm}
\usepackage{algorithmic}
\usepackage{placeins}
\usepackage{dsfont}
\usepackage{url}
\usepackage{graphics}
\usepackage{adjustbox}
\usepackage{todonotes}
\usepackage{soul}
\usepackage{dsfont}
\usepackage{algorithm}
\usepackage{algorithmic}
\newcommand{\junk}[1]{}
\newtheorem{theorem}{Theorem}[section]
\newtheorem{corollary}{Corollary}[theorem]
\newtheorem{lemma}[theorem]{Lemma}
\newtheorem{proposition}[theorem]{Proposition}
\begin{document}
\title[Higher Accuracy Modular Data Assimilation for the Navier-Sto]{Higher Accuracy Modular Data Assimilation for the Navier-Stokes Equations}
\author{Troy Yang}
\date{}
\maketitle
\noindent\emph{Source and licence.} Higher Accuracy Modular Data Assimilation for the Navier-Stokes Equations. Version arXiv:2606.19508v2 (CC BY 4.0). This material is distributed under the Creative Commons Attribution 4.0 International licence, http://creativecommons.org/licenses/by/4.0/, as recorded in the arXiv OAI licence metadata for this exact version and shown on the abstract page https://arxiv.org/abs/2606.19508v2. Full rights notice: licenses/arxiv-2606.19508-CC-BY-4.0.txt.\par\smallskip
\noindent\textit{Editorial scope.} Counted unit: the original \texttt{theorem} environment \texttt{-} at source lines 594--611 together with its complete proof at lines 613--661; the delivered document keeps source lines 355--662 so that every referenced label and every citation resolves inside it. Native editable source is the paper's own concatenated TeX; the AMS wrapper is editorial formatting only. No publisher class, upstream script or unrelated artwork is redistributed.
\section{Error Analysis}
In this section, we prove two error estimates for the 2-step nudging with the analysis step in (1.5). One estimate imposes no condition on the time step size, while another estimate imposes a small time step size condition with restrictions on $\chi$. Before we prove these estimates, we introduce a few lemmas necessary for proving the error estimates.


\begin{proposition}
    Consider the analysis step in (4.3). Then the following identities holds:
    \begin{equation}
        ||e^{n+2}||^2+||e^{n+2}-\tilde{e}^{n+2}||^2+\frac{4}{3}\Delta t\chi ||I_He^{n+2}||^2= ||\tilde{e}^{n+2}||^2
    \end{equation}
    and
    \begin{equation}
        ||I_He^{n+2}||^2+||I_H(e^{n+2}-\tilde{e}^{n+2})||^2+\frac{4}{3}\Delta t\chi ||I_He^{n+2}||^2= ||I_H\tilde{e}^{n+2}||^2.
    \end{equation}
\end{proposition}

\begin{proof}
    At time $t^{n+2}$, the true solution $u$ satisfies
    \begin{align}
        \frac{3u^{n+2}-3u^{n+2}}{2\Delta t} -\chi I_H(u(t^{n+2})-u^{n+2})=0.
    \end{align}
    Subtract (1.5) from (5.3) and take inner product with $e^{n+2}$ gives
    \begin{align}
        \Bigg(\frac{3e^{n+2}-3\Tilde{e}^{n+2}}{2\Delta t}, e^{n+2}\Bigg)
        &+\chi (I_He^{n+2}, e^{n+2})=0.
    \end{align}
    For the first term on LHS of (5.4), we apply polarization identity. For the second term, we apply $I_H$ to right side of the inner product. Then
    \begin{align}
        \frac{3}{2\Delta t}\Bigg(\frac{1}{2}||e^{n+2}||^2-\frac{1}{2}||\tilde{e}^{n+2}||^2+\frac{1}{2}||e^{n+2}-\tilde{e}^{n+2}||^2\Bigg)+\chi||I_H e^{n+2}||^2=0.
    \end{align}
    Rearranging terms completes the proof of (5.1). To prove (5.2), take inner product with $I_He^{n+2}$ and apply $I_H$ to the left side of the first term. Then
    \begin{align}
        \frac{3}{2\Delta t}(I_H(e^{n+2}-\Tilde{e}^{n+2}), I_He^{n+2})
        &+\chi (I_He^{n+2}, I_He^{n+2})=0.
    \end{align}
    We apply polarization identity to the first term. We have
    \begin{align}
        \frac{3}{2\Delta t}\Bigg(\frac{1}{2}||I_He^{n+2}||^2-\frac{1}{2}||I_H\tilde{e}^{n+2}||^2+\frac{1}{2}||I_H(e^{n+2}-\tilde{e}^{n+2})||^2\Bigg)+\chi||I_H e^{n+2}||^2=0.
    \end{align}
    Rearranging terms completes the proof of (5.2).
\end{proof}

Remark: From (5.1) of Proposition 5.1, we know $||e^{n+2}||^2=||\tilde{e}^{n+2}||^2$ only if $||e^{n+2}-\tilde{e}^{n+2}||^2=||I_He^{n+2}||^2=0$. With (5.2) of Proposition 5.1, we know $||I_He^{n+2}||^2=||I_H\tilde{e}^{n+2}||^2$ only when $||I_H(e^{n+2}-\tilde{e}^{n+2})||^2=0$ and $(1+\frac{4}{3}\Delta t \chi)||I_He^{n+2}||^2=||I_H\tilde{e}^{n+2}||^2$.

\begin{corollary}
    Consider the analysis step in (4.3). Then the following identity holds:
    \begin{align}
        ||e^{n+2}||^2< ||\tilde{e}^{n+2}||^2
    \end{align} with equality when $e^{n+2}=\tilde{e}^{n+2}$ and $I_He^{n+2}=0.$
\end{corollary}

\begin{proof}
    Take (5.1). Since all terms are positive and there are more terms on LHS, we have $||e^{n+2}||^2<||\tilde{e}^{n+2}||^2$. To obtain equality, we need $||e^{n+2}-\tilde{e}^{n+2}||^2=0$ and $||I_He^{n+2}||^2=0$. This only happens when $e^{n+2}=\tilde{e}^{n+2}$ and $I_He^{n+2}=0$.
\end{proof}

\begin{corollary}
    Consider the analysis step in (4.3). Then the following identity holds:
    \begin{align}
        ||I_He^{n+2}||^2< ||I_H\tilde{e}^{n+2}||^2
    \end{align}
    with equality when $I_He^{n+2}=I_H\tilde{e}^{n+2}=0$.
\end{corollary}

\begin{proof}
    Take (5.2). Since all terms are positive and there are more terms on LHS, we have $||I_He^{n+2}||^2<||I_H\tilde{e}^{n+2}||^2$. To obtain equality, we need $||I_H(e^{n+2}-\tilde{e}^{n+2})||^2=0$ and $||I_He^{n+2}||^2=0$. This only happens when $I_He^{n+2}=I_H\tilde{e}^{n+2}=0$.
\end{proof}

\begin{lemma}
    Consider the analysis step in (4.3). If $I_H$ is $L^2$ orthogonal projection, then the following identity holds:
    \begin{equation}
        ||I_H\tilde{e}^{n+2}||^2=\Bigg(1+\frac{2}{3}\Delta t\chi\Bigg)^2||I_He^{n+2}||^2.
    \end{equation}
\end{lemma}

\begin{proof}
    At time $t^{n+2}$, the error equation is
    \begin{align}
        \Bigg(\frac{3e^{n+2}-3\Tilde{e}^{n+2}}{2\Delta t},v_h\Bigg) +\chi (I_He^{n+2}, v_h)=0.
    \end{align}
    We choose $v_h=I_He^{n+2}$. Then
    \begin{align}
        \frac{3}{2\Delta t}(e^{n+2}-\tilde{e}^{n+2}, I_He^{n+2})+\chi||I_He^{n+2}||^2=0.
    \end{align}
    We apply $I_H$ projection to the first term and rearrange terms to get
    \begin{align}
        (I_H\tilde{e}^{n+2}, I_He)=\Bigg(1+\frac{2}{3}\Delta t\chi \Bigg)||I_He^{n+2}||^2.
    \end{align}
    Now we choose $v_h=I_H\tilde{e}^{n+2}$. Then
    \begin{align}
        \frac{3}{2\Delta t}(e^{n+2}-\tilde{e}^{n+2}, I_H\tilde{e}^{n+2})+\chi(I_He^{n+2}, I_H\tilde{e}^{n+2})=0.
    \end{align}
    We again apply $I_H$ projection to the first term and rearrange terms to get
    \begin{align}
        (I_He^{n+2}, I_H\tilde{e}^{n+2})=\frac{1}{1+\frac{2}{3}\Delta t \chi}||I_H\tilde{e}^{n+2}||^2.
    \end{align}
    Due to symmetry, we can combine (5.13) and (5.15) to yield
    \begin{align}
        \Bigg(1+\frac{2}{3}\Delta t \chi\Bigg)||I_He^{n+2}||^2=\frac{1}{1+\frac{2}{3}\Delta t\chi}||I_H\tilde{e}^{n+2}||^2.
    \end{align}
    Multiplying through by $\Bigg(1+\frac{2}{3}\Delta t\chi\Bigg)$ completes the proof.
\end{proof}

The following lemma is crucial for handling the time discretization term in the error analysis.

\begin{lemma}
    Consider the analysis step in (4.3). Then the following identity holds:
    \begin{align}
        &\Bigg(\frac{3e^{n+2}-4e^{n+1}+e^n}{2\Delta t}, \Tilde{e}^{n+2}-e^{n+2}\Bigg) \notag \\
        =&\frac{1}{6}\chi \Bigg[ ||I_He^{n+2}||^2-||I_He^{n+1}||^2 + ||2I_He^{n+2}-I_He^{n+1}||^2 - ||2I_He^{n+1}-I_He^n||^2 \\
        &+ ||I_He^{n+2}-2I_He^{n+1}+I_He^n||^2\Bigg]. \notag
    \end{align}
\end{lemma}

\begin{proof}
    At time $t^{n+2}$, the true solution $u$ satisfies
    \begin{align}
        \frac{3u^{n+2}-3u^{n+2}}{2\Delta t} -\chi I_H(u(t^{n+2})-u^{n+2})=0.
    \end{align}
    Subtract (1.5) from (5.18) and take inner product with $\frac{3e^{n+2}-4e^{n+1}+e^n}{3}$ gives
    \begin{align}
        \Bigg(\frac{3e^{n+2}-3\Tilde{e}^{n+2}}{2\Delta t}, \frac{3e^{n+2}-4e^{n+1}+e^n}{3}\Bigg)
        &+\chi \Bigg(I_He^{n+2}, \frac{3e^{n+2}-4e^{n+1}+e^n}{3}\Bigg)=0. \notag
    \end{align}
    For the second term, we apply $I_H$ to right side of the inner product and use Lemma 2.1. Rearranging terms completes the proof.
\end{proof}

We now present the first error estimate in Theorem 5.4.

\begin{theorem}
    Consider forecast step in (1.3) and analysis step in (1.5). Suppose $I_H$ is an $L^2$ projection. Assume 
    \begin{align}
        \frac{31}{8}\nu-4\frac{U}{l}C_P^2L^2>0 \text{, equivalently, } Re=\frac{UL}{\nu}<\frac{31}{32}\frac{l}{C_P^2L}. \notag
    \end{align} Then the following holds:
    \begin{align}
        ||e^N||^2
        &+||2e^N-e^{N-1}||^2+\frac{2}{3}\Delta t\chi ||I_He^N||^2+\frac{2}{3}\Delta t\chi ||2I_He^N-I_He^{N-1}||^2 \notag \\
        &+\sum_{n=0}^{N-2}||e^{n+2}-2e^{n+1}+e^n||^2 + \frac{2}{3}\Delta t\chi \sum_{n=0}^{N-2} ||I_H(e^{n+2}-2e^{n+1}+e^n)||^2 \notag \\
        &+ 6 \sum_{n=0}^{N-2} ||\tilde{e}^{n+2}-e^{n+2}||^2 +\Delta t\Bigg(\frac{31}{8}\nu -4\frac{U}{l}C_P^2L^2\Bigg) \sum_{n=0}^{N-2}||\nabla \tilde{e}^{n+2}||^2 \\
        \leq \frac{8}{\nu}&C\Delta t^4\int_{t^0}^{t^{N}}||u_{ttt}||_{-1}^2 + ||e^1||^2+||2e^1-e^0||^2 \notag \\
        &+\frac{2}{3}\Delta t\chi ||I_He^1||^2+\frac{2}{3}\Delta t\chi ||2I_He^1-I_He^0||^2. \notag
    \end{align}
\end{theorem}

\begin{proof}
    For the forecast step, at time $t^{n+2}$, the true solution $u,p$ satisfies
    \begin{align}
        \frac{3u^{n+2}-4u^{n+1}+u^n}{2\Delta t}+u^{n+2} \cdot \nabla u^{n+2} - \nu \Delta u^{n+2} + \nabla p^{n+2}=f^{n+2}+\rho^{n+2},
    \end{align}
    where $\rho^{n+2}=\frac{3u^{n+2}-4u^{n+1}+u^n}{2\Delta t}-u_t(t^{n+2})$ is the consistency error term.
    
    Subtracting (1.3) from (5.20) and taking inner product with $\Tilde{e}^{n+2}$ yields
    \begin{align}
        \Bigg(\frac{3\Tilde{e}^{n+2}-4e^{n+1}+e^n}{2\Delta t}, \Tilde{e}^{n+2}\Bigg)
        &+(u^{n+2} \cdot \nabla u^{n+2}- \tilde{v}^{n+2} \cdot \nabla \Tilde{v}^{n+2}, \Tilde{e}^{n+2}) \notag \\
        &+ \nu ||\nabla \Tilde{e}^{n+2}||^2=(\rho^{n+2}, \Tilde{e}^{n+2}).
    \end{align}
    We rewrite the first term exactly like in (4.12). Now, the first term is split into three terms. For the split terms, we apply Lemma 2.1 to the first term, polarization identity to the second term, and (5.17) of Lemma 5.3 to the third term. Then, we have
    \begin{align}
        &\frac{1}{4\Delta t}\Bigg[||e^{n+2}||^2-||e^{n+1}||^2+||2e^{n+2}-e^{n+1}||^2-||2e^{n+1}-e^n||^2+||e^{n+2}-2e^{n+1}+e^n||^2\Bigg] \notag   \\
        &+\frac{1}{6}\chi \Bigg[ ||I_He^{n+2}||^2 - ||I_He^{n+1}||^2 + ||2I_He^{n+2}-I_He^{n+1}||^2 - ||2I_He^{n+1}-I_He^n||^2  \notag \\
        &+||I_He^{n+2}-2I_He^{n+1}+I_He^n||^2\Bigg] \\
        &+\frac{3}{4\Delta t}\Bigg[||\Tilde{e}^{n+2}||^2-||e^{n+2}||^2 + ||\Tilde{e}^{n+2}-e^{n+2}||^2 \Bigg] \notag \\
        &+(u^{n+2} \cdot \nabla u^{n+1}- \tilde{v}^{n+2} \cdot \nabla \Tilde{v}^{n+2}, \Tilde{e}^{n+2}) +\nu ||\nabla \Tilde{e}^{n+2}||^2 = (\rho^{n+2},\Tilde{e}^{n+2}). \notag
    \end{align}
    We handle the nonlinear term by adding and subtracting $\tilde{v}^{n+2}\cdot \nabla u^{n+2}$ as follows
    \begin{align}
        &(u^{n+2} \cdot \nabla u^{n+2}-\tilde{v}^{n+2} \cdot \nabla \Tilde{v}^{n+2}, \Tilde{e}^{n+2}) \notag \\
        =&(u^{n+2} \cdot \nabla u^{n+2} -\tilde{v}^{n+2} \cdot \nabla u^{n+2} + \tilde{v}^{n+2}\cdot \nabla u^{n+2} + \tilde{v}^{n+2}\cdot \nabla \tilde{v}^{n+2}, \Tilde{e}^{n+2}) \notag \\
        =&(\tilde{e}^{n+2}\cdot \nabla u^{n+2}, \Tilde{e}^{n+2})+(\tilde{v}^{n+2}\cdot \nabla \Tilde{e}^{n+2}, \Tilde{e}^{n+2}) \notag \\
        =&(\tilde{e}^{n+2}\cdot \nabla u^{n+2}, \Tilde{e}^{n+2}).    \end{align}
    We bound (5.23) using H\"{o}lder's inequality as follows
    \begin{align}
        (\tilde{e}^{n+2}\cdot \nabla u^{n+2}, \Tilde{e}^{n+2}) 
        \leq ||\tilde{e}^{n+2}||\cdot ||\nabla u^{n+2}||_\infty \cdot ||\Tilde{e}^{n+2}|| &\leq \frac{U}{l}||\tilde{e}^{n+2}||^2 \notag \\
        &\leq \frac{U}{l}C_P^2L^2 ||\nabla \tilde{e}^{n+2}||^2,
    \end{align}
    where we use Poincar\'e for the last inequality. For the consistency term, we bound using H\"{o}lder's and Young's inequalities as follows
    \begin{align}
        (\rho^{n+2}, \Tilde{e}^{n+2})\leq ||\rho^{n+2}||_{-1} \cdot ||\nabla \Tilde{e}^{n+2}||\leq \frac{2}{\nu}||\rho^{n+2}||_{-1}^2+\frac{\nu}{8}||\nabla \Tilde{e}^{n+2}||^2.
    \end{align}
    Using Lemma 2.2, 
    \begin{align}
        ||\rho^{n+2}||_{-1}^2\leq C\Delta t^3\int_{t^n}^{t^{n+2}}||u_{ttt}||_{-1}^2dt.
    \end{align}
    For the analysis step, at time $t^{n+2}$, the true solution $u$ satisifies
    \begin{align}
        \frac{3u^{n+2}-3u^{n+2}}{2\Delta t}-\chi I_H(u(t^{n+2})-u^{n+2})=0.
    \end{align}
    Subtracting (1.5) from (5.18) and taking inner product with $e^{n+2}$ gives
    \begin{align}
        \Bigg(\frac{3e^{n+2}-3\Tilde{e}^{n+2}}{2\Delta t}, e^{n+2}\Bigg)+\chi (I_He^{n+2},e^{n+2})=0.
    \end{align}
    We apply the polarization identity to the first term. For the second term, we apply $I_H$ to the right side of the inner product. Rearranging terms, we have 
    \begin{align}
        ||e^{n+2}||^2=||\tilde{e}^{n+2}||^2-||e^{n+2}-\tilde{e}^{n+2}||^2-\frac{4}{3}\Delta t\chi ||I_He^{n+2}||^2.
    \end{align}
    We use (5.29) to replace $||e^{n+2}||^2$ in (5.22) and combine the resulting equality with (5.24) and (5.26). Rearranging terms, we have
    \begin{align}
        &\Bigg[||e^{n+2}||^2-||e^{n+1}||^2+||2e^{n+2}-e^{n+1}||^2-||2e^{n+1}-e^n||^2+||e^{n+2}-2e^{n+1}+e^n||^2\Bigg] \notag  \\
        &+\frac{2}{3}\Delta t\chi \Bigg[ ||I_He^{n+2}||^2 - ||I_He^{n+1}||^2 + ||2I_He^{n+2}-I_He^{n+1}||^2 - ||2I_He^{n+1}-I_He^n||^2 \notag \\
        &+||I_He^{n+2}-2I_He^{n+1}+I_He^n||^2\Bigg]  \\
        &+6||\Tilde{e}^{n+2}-e^{n+2}||^2 +4\Delta t\chi ||I_He^{n+2}||^2 +\Delta t\Bigg(\frac{31}{8}\nu -4\frac{U}{l}C_P^2L^2\Bigg) ||\nabla \Tilde{e}^{n+2}||^2 \notag \\
        &\leq \frac{8}{\nu}C\Delta t^4\int_{t^n}^{t^{n+2}}||u_{ttt}||_{-1}^2dt. \notag
    \end{align}
    Taking the sum from $n = 0$ to $N-2$
    completes the proof.
\end{proof}

\begin{corollary}
    Consider forecast step in (1.3) and analysis step in (1.5). Suppose $I_H$ is an $L^2$ projection. Assume 
    \begin{align}
        \frac{31}{8}\nu-4\frac{U}{l}C_P^2L^2>0  \text{, equivalently, } Re=\frac{UL}{\nu}<\frac{31}{32}\frac{l}{C_P^2L}. \notag
    \end{align} Then the following holds:
    \begin{align}
        ||e^N||^2\leq C&\Delta t^4\int_{t^0}^{t^{N}}||u_{ttt}||_{-1}^2 + ||e^1||^2+||2e^1-e^0||^2 \notag \\
        &+\frac{2}{3}\Delta t\chi ||I_He^1||^2+\frac{2}{3}\Delta t\chi ||2I_He^1-I_He^0||^2.
    \end{align}
\end{corollary}

\begin{proof}
    We begin with (5.19) and assume $\frac{31}{8}\nu-4\frac{U}{l}C_P^2>0$. Since all LHS terms are nonnegative, we can drop all but $||e^N||^2$. On the RHS, we have the integral term and initial errors, bounding the error at time $N$.
\end{proof}

\begin{corollary}
    Consider forecast step in (1.3) and analysis step in (1.5). Suppose $I_H$ is an $L^2$ projection. Assume 
    \begin{align}
        \frac{31}{8}\nu-4\frac{U}{l}C_P^2L^2>0 \text{, equivalently, } Re=\frac{UL}{\nu}<\frac{31}{32}\frac{l}{C_P^2L}. \notag
    \end{align} Then the following holds:
    \begin{align}
        ||e^N||=\mathcal{O}(\Delta t^2).
    \end{align}
\end{corollary}

\begin{proof}
    We begin with (5.19). We take the square root of both sides of the inequality. Since all LHS terms are nonnegative, drop all but $||e^N||^2$. On the RHS, the dominant term is the integral term with coefficient $\Delta t^2$. Thus, the error at time $N$ is second-order accurate in time. 
\end{proof}
From the two corollaries above, we know the error does not grow in time, the numerical solution remains uniformly bounded for all $N$, and the method converges second-order in time.

Now we consider a different strategy in proving the error estimate and present the small timestep estimate in Theorem 5.5.


\begin{theorem}
    Consider forecast step in (1.3) and analysis step in (1.5). Suppose $I_H$ is an $L^2$ projection. Assume 
    \begin{align}
        \nu-\frac{8}{7}\frac{U}{l}C^2H^2>0 \text{, equivalently, } Re=\frac{UL}{\nu}<\frac{7}{8}\frac{lL}{C^2H^2}. \notag
    \end{align}
    Define $T^*=\frac{L}{U}$. Assume
    \begin{align}
    \frac{\Delta t}{T^*} \leq \frac{3}{8}\frac{l}{L}. \notag
    \end{align}
    If $\chi\in [\chi_1,\chi_2]\cap(0,\infty)$, where $\chi_1,\chi_2$ are the roots of $\chi-\frac{U}{l}(1+\frac{2}{3}\Delta t\chi)^2$, then the following holds:
    \begin{align}
        ||e^N||^2
        &+||2e^N-e^{N-1}||^2+\frac{2}{3}\Delta t\chi ||I_He^N||^2+\frac{2}{3}\Delta t\chi ||2I_He^N-I_He^{N-1}||^2  \notag \\
        +\sum_{n=0}^{N-2}&||e^{n+2}-2e^{n+1}+e^n||^2 + \frac{2}{3}\Delta t\chi \sum_{n=0}^{N-2} ||I_H(e^{n+2}-2e^{n+1}+e^n)||^2 + 6 \sum_{n=0}^{N-2} ||\tilde{e}^{n+2}-e^{n+2}||^2 \notag \\
        +\Delta t&\Bigg(\frac{7}{2}\nu -4\frac{U}{l}C^2H^2\Bigg) \sum_{n=0}^{N-2}||\nabla \tilde{e}^{n+2}||^2+4\Delta t\Bigg(\chi - \frac{U}{l}(1+\frac{2}{3}\Delta t\chi)^2\Bigg)\sum_{n=0}^{N-2} ||I_He^{n+2}||^2 \\
        \leq \frac{8}{\nu}&C\Delta t^4\int_{t^0}^{t^{N}}||u_{ttt}||_{-1}^2 + ||e^1||^2+||2e^1-e^0||^2+\frac{2}{3}\Delta t\chi ||I_He^1||^2+\frac{2}{3}\Delta t\chi ||2I_He^1-I_He^0||^2. \notag
    \end{align}
\end{theorem}

\begin{proof}
    The proof starts identical to the proof of Theorem 5.4 but deviates by a different bound for (5.24).
    We bound (5.24) using H\"{o}lder's inequality as follows
    \begin{align}
        (\tilde{e}^{n+2}\cdot \nabla u^{n+2}, \Tilde{e}^{n+2}) 
        &\leq ||\tilde{e}^{n+2}||\cdot ||\nabla u^{n+2}||_\infty \cdot ||\Tilde{e}^{n+2}|| \notag \\
        &\leq \frac{U}{l}||\tilde{e}^{n+2}||^2, \text{ where } \frac{U}{l}=\max_{\Omega \times (0,\infty)}||\nabla u(x,t)||, U= \max_{\Omega \times (0,\infty)}|| u(x,t)||\notag \\
        &\leq \frac{U}{l}\Bigg[||I_H\tilde{e}^{n+2}||^2+ ||(I-I_H)\tilde{e}^{n+2}||^2\Bigg] \notag \\
        &\leq \frac{U}{l}||I_H\tilde{e}^{n+2}||^2+ \frac{U}{l}C^2H^2||\nabla \tilde{e}^{n+2}||^2,
    \end{align}
    where the last inequality follows from (2.1).
    For the consistency term, we bound using H\"{o}lder's and Young's inequalities as follows
    \begin{align}
        (\rho^{n+2}, \Tilde{e}^{n+2})\leq ||\rho^{n+2}||_{-1} \cdot ||\nabla \Tilde{e}^{n+2}||\leq \frac{2}{\nu}||\rho^{n+2}||_{-1}^2+\frac{\nu}{8}||\nabla \Tilde{e}^{n+2}||^2.
    \end{align}
    Using Lemma 2.2, 
    \begin{align}
        ||\rho^{n+2}||_{-1}^2\leq C\Delta t^3\int_{t^n}^{t^{n+2}}||u_{ttt}||_{-1}^2dt.
    \end{align}
    For the analysis step, at time $t^{n+2}$, the true solution $u$ satisifies
    \begin{align}
        \frac{3u^{n+2}-3u^{n+2}}{2\Delta t}-\chi I_H(u(t^{n+2})-u^{n+2})=0.
    \end{align}
    Subtracting (1.5) from (5.37) and taking inner product with $e^{n+2}$ gives
    \begin{align}
        \Bigg(\frac{3e^{n+2}-3\Tilde{e}^{n+2}}{2\Delta t}, e^{n+2}\Bigg)+\chi (I_He^{n+2},e^{n+2})=0.
    \end{align}
    We apply the polarization identity to the first term. For the second term, we apply $I_H$ to the right side of the inner product. Rearranging terms, we have 
    \begin{align}
        ||e^{n+2}||^2=||\tilde{e}^{n+2}||^2-||e^{n+2}-\tilde{e}^{n+2}||^2-\frac{4}{3}\Delta t\chi ||I_He^{n+2}||^2.
    \end{align}
    We use (5.39) to replace $||e^{n+2}||^2$ in (5.22) and combine the resulting equality with (5.34) and (5.36). Rearranging terms, we have
    \begin{align}
        \Bigg[||e^{n+2}||^2-||e^{n+1}||^2+&||2e^{n+2}-e^{n+1}||^2-||2e^{n+1}-e^n||^2+||e^{n+2}-2e^{n+1}+e^n||^2\Bigg] \notag  \\
        +\frac{2}{3}\Delta t\chi \Bigg[||I_He^{n+2}&||^2 - ||I_He^{n+1}||^2 + ||2I_He^{n+2}-I_He^{n+1}||^2 - ||2I_He^{n+1}-I_He^n||^2 \notag \\
        +||&I_He^{n+2}-2I_He^{n+1}+I_He^n||^2\Bigg]  \\
        +6||\Tilde{e}^{n+2}-e^{n+2}&||^2 +4\Delta t\chi ||I_He^{n+2}||^2 +\Delta t\Bigg(\frac{7}{2}\nu -4\frac{U}{l}C^2H^2\Bigg) ||\nabla \Tilde{e}^{n+2}||^2 \notag \\
        \leq \frac{8}{\nu}C\Delta t^4\int_{t^n}^{t^{n+2}}&||u_{ttt}||_{-1}^2dt + 4\Delta t\frac{U}{l}||I_H\tilde{e}^{n+2}||^2. \notag
    \end{align}
    Applying Lemma 5.2 to $||I_H\tilde{e}^{n+2}||^2$ and rearranging terms give
    \begin{align}
        \Bigg[||e^{n+2}||^2-||e^{n+1}||^2+&||2e^{n+2}-e^{n+1}||^2-||2e^{n+1}-e^n||^2+||e^{n+2}-2e^{n+1}+e^n||^2\Bigg] \notag  \\
        +\frac{2}{3}\Delta t\chi \Bigg[ ||I_He^{n+2}||^2& - ||I_He^{n+1}||^2 + ||2I_He^{n+2}-I_He^{n+1}||^2 - ||2I_He^{n+1}-I_He^n||^2 \notag \\
        &+||I_He^{n+2}-2I_He^{n+1}+I_He^n||^2\Bigg]  \\
        +6||\Tilde{e}^{n+2}-e^{n+2}||^2 &+4\Delta t\Bigg(\chi-\frac{U}{l}(1+\frac{2}{3}\Delta t \chi)^2\Bigg)||I_He^{n+2}||^2 \notag \\
        +\Delta t\Bigg(\frac{7}{2}\nu -4\frac{U}{l}C^2&H^2\Bigg) ||\nabla \Tilde{e}^{n+2}||^2 \leq \frac{8}{\nu}C\Delta t^4\int_{t^n}^{t^{n+2}}||u_{ttt}||_{-1}^2dt. \notag
    \end{align}
    Taking the sum from $n = 0$ to $N-2$ completes the proof.
\end{proof}


\end{document}
